\section{The specified lattice Hamiltonian}
\label{notes-model}

We use the spatial Wilson Hamiltonian of the compact $U(1)$
framework~\cite{KogutSusskind,Zache}, with continuous time and untruncated
electric rotors. Let $X=(\Z/n\Z)^3$, $n\ge3$, $N=n^3$, and let
$\mathcal E$ be its $3N$ positively oriented links.
The lattice spacing is $a>0$, coupling $g^2>0$, bare mass $m\ge0$,
and Wilson parameter $r\ge0$. The Hilbert space is
\[
 \mathcal H=L^2(U(1)^{\mathcal E},dU)\otimes
 \mathcal F_a(\ell^2(X)\otimes\C^4).
\]
All Haar measures are normalized. A link variable is
$U_\ell=e^{i\varphi_\ell}$ and $E_\ell=-i\partial_{\varphi_\ell}$
has integer spectrum. With one plaquette for each $x,j<k$, set
\[
 H_E=\frac{g^2}{2a}\sum_\ell E_\ell^2,\qquad
 H_B=\frac1{g^2a}\sum_p(1-\operatorname{Re}U_p).
\]
The four Hermitian Dirac matrices $\alpha_j,\beta$ anticommute
pairwise and square to $I$. Spatial operators commute with them.
The negatively charged fermion convention is
\begin{align}
 (T_j(U)f)(x)&=\overline{U_{xj}}f(x+e_j),&
 v_j(U)&=\frac{T_j(U)-T_j(U)^*}{2ia},\label{lat-hopping}\\
 M(U)&=m+\frac r{2a}\sum_j(2-T_j(U)-T_j(U)^*),&&\label{lat-wilson-mass}\\
 h_W(U)&=\sum_j\alpha_jv_j(U)+\beta M(U),&
 H&=H_E+H_B+d\Gamma(h_W(U)).\label{lat-hamiltonian}
\end{align}
Each $T_j$ is unitary, so
$M(U)\ge mI$ and $\|h_W(U)\|\le h_*:=m+(3+6r)/a$.
The entire raw bilinear $d\Gamma(h_W(U))$ is retained.
Changing to electron--positron variables retains its link-dependent
scalar sea term. Dropping that term defines a different Hamiltonian.

Let $q_x=\sum_{s=1}^4c_{xs}^*c_{xs}-2$, so that
$q_x\in\{-2,-1,0,1,2\}$ and physical fermion charge is $-gq_x$.
The strongly commuting Gauss generators are
\[
 G_x=\sum_j(E_{xj}-E_{x-e_j,j})+q_x,\qquad
 \mathcal H_\nu=\{\Psi:G_x\Psi=\nu_x\Psi\ \text{for all }x\}.
\]
Sources are integer valued. A sector is nonempty precisely when
$|\sum_x\nu_x|\le2N$: necessity follows by summing Gauss's law;
for sufficiency choose occupations with that total and route the
remaining zero-sum integer function along a spanning tree.
For nuclei $\nu_x=\sum_i Z_i\mathbf1_{\{x=R_i\}}$.
Write $E_\nu=\inf\sigma(H|_{\mathcal H_\nu})$ and $E_0$ for the
physical interacting vacuum. There is no extra global-flux restriction.

\begin{lemma}[Realization and elementary energy bounds]\label{lat-realization}
The Hamiltonian is self-adjoint on $D(H_E)$, has compact resolvent,
and has trace-class heat semigroup for every positive time. These
properties hold in every admissible Gauss sector. Moreover,
\[
 H\ge-2Nh_*,\qquad
 E_0\le-2N(m+3r/a)+\frac{3N}{g^2a}.
\]
\end{lemma}
\begin{proof}
The remaining terms are bounded matrix multipliers, so bounded
perturbation of the torus Laplacian proves all operator assertions,
including the heat-trace comparison. The matrix
$C_5=\alpha_1\alpha_2\alpha_3\beta$ anticommutes with $h_W$,
whose eigenvalues therefore come in opposite pairs.
The least fermionic Fock energy at fixed $U$ is
$-\frac12\Tr|h_W(U)|\ge-2Nh_*$. For the upper bound use a constant
link wavefunction and fill the two negative-$\beta$ spinors at every
site. It is a physical vacuum trial state with zero electric energy,
zero hopping expectation, and the displayed mass and magnetic energies.
\end{proof}

\begin{theorem}[Source comparison recalled from Paper A]\label{lat-main}
For every admissible integer source and the above parameters,
\[
 E_\nu\ge E_0+m\left|\sum_x\nu_x\right|.
\]
This includes source electric self-energy. With the individual
periodic lattice Coulomb diagonals subtracted, positive nuclei
satisfy a lower bound $-C_aK$ relative to the same vacuum,
uniformly in volume and positions.
\end{theorem}
The proof, including the finite combinatorial factor at positive
temperature, is in Paper A~\cite{PaperA}. Its inputs are the centered
fermion determinant, $C_5$ pseudo-unitarity, $\beta$-form decrease,
the positive electric heat kernel, and exact local Gauss projection.
Sections below prove the separate static comparisons and obstructions;
they do not infer a neutral pair gap from this net-charge theorem.

For the self-energy statement, the convention is
$S_n=\frac{g^2}{2a}G_n(0)\sum_iZ_i^2$ with
$G_n(0)=n^{-3}\sum_{k\ne0}[4\sum_j\sin^2(\pi k_j/n)]^{-1}$.
The elementary uniform bound $G_n(0)\le13/16$ gives
$E_\nu-S_n-E_0\ge m\sum_iZ_i-(13g^2/(32a))Z_{\max}^2K$.
This specifies the comparison used for the clover variant; no
configuration-dependent subtraction is introduced.
